\exercisegroup

\begin{enumerate}
\def\labelenumi{\arabic{enumi})}
\item
  设 X、Y 是拓扑空间，且 \(f: X \to Y\) 是连续映射。设 \((T, p)\) 是 Y-空间。X-空间上的典范下降数据 \(\tau\) 在 \(Z = X \times_{Y} T\) 上是有效的（IV，第 385 页，命题 1），典范映射 \(Z/R_{\tau} \to T\) 是单射且连续的。然而，它不一定是满射，也不一定严格。
\item
  设 Y 是由各空间 \(X_{i}\) 粘合得到的拓扑空间。令 X 为族 \((\mathrm{X}_{i})\) 的不交并所成的拓扑空间，令 \(f: X \to Y\) 为典范映射（TG，第 I 章，第 16 页）。对每个 i，置 \(Y_{i} = f(\mathrm{X}_{i})\)，并令 \(f_{i}: X_{i} \to Y_{i}\) 为 f 通过取子空间诱导的映射。映射 \(f_{i}\) 是双射且连续的。若映射 \(f_{i}\) 不是同胚，则相对于 f 的 X 上的典范下降数据不是有效的。例如参见 TG，第 I 章，第 94 页，练习 15。
\item
  令 B 为点 \((x,y)\) 在正方形 \([0,1]^{2}\) 中且满足 x=0、或 x=1、或 y=0、或 y=1、或存在 \(n\in N\) 使 nx=1 所成的集合。令 \(B_{1}\) 和 \(B_{2}\) 分别为 B 的子集，由 B 中的点 \((x,y)\) 组成，其中分别满足 \(y\leqslant\frac{1}{2}\) 和 \(y\geqslant\frac{1}{2}\)。
\end{enumerate}

构造一个不是覆叠的 B-空间 \((\mathrm{E}, p)\)，但使得 \(B_{1}\)-空间 \((\mathrm{E}_{\mathrm{B}_{1}}, p_{\mathrm{B}_{1}})\) 和 \(B_{2}\)-空间 \((\mathrm{E}_{\mathrm{B}_{2}}, p_{\mathrm{B}_{2}})\) 都是覆叠。

\begin{enumerate}
\def\labelenumi{\arabic{enumi})}
\setcounter{enumi}{3}
\tightlist
\item
  若拓扑空间的每个点都有由道路单连通邻域组成的邻域基，则称该空间局部道路单连通。
\end{enumerate}

\begin{enumerate}
\def\labelenumi{\alph{enumi})}
\item
  举出一个道路单连通但不是局部道路单连通的拓扑空间的例子。
\item
  设 X 是拓扑空间，G 是在 X 上适当作用的离散群。若 X 局部道路单连通，则 X/G 也是如此。
\end{enumerate}

\begin{enumerate}
\def\labelenumi{\arabic{enumi})}
\setcounter{enumi}{4}
\item
  设 G 在可解圈且完全正则的拓扑空间 X 上适当作用。置 Y = X/G，并记 \(f: X \to Y\) 为典范映射。证明典范群胚态射 \(\varpi(X)/G \to \text{Coeg}(f)\) 是同构。（使用 LIE，第 IX 章，§9，练习 17，改写 IV，第 403 页定理 3 的证明。）
\item
 设 X 是道路连通的拓扑空间，G 是连续作用于 X 上的群，且 x 是 X 中满足 \(g \cdot x = x\)（对所有 \(g \in G\)）的点。记 p 为从 X 到 X/G 的典范映射。设 H 是群胚，且 \(\varphi: \varpi(X) \to H\) 是群胚态射。假设对于每条道路 \(c\)，它在 \(\mathrm{X}\) 上，并且对于每个 \(g\in \mathrm{G}\)，都有 \(\varphi ([g\cdot c]) = \varphi ([c])\)，并且对于每个圈 \(c\in \Omega_x(\mathrm{X})\)，都有 \(\varphi ([c]) = e_{\varphi (x)}\)。
\end{enumerate}

设 \(c_{1}\) 和 \(c_{2}\) 是 X 中的道路，使得 \(p \circ c_{1}\) 与 \(p \circ c_{2}\) 具有相同的起点和相同的终点。证明 \(\varphi([c_{1}]) = \varphi([c_{2}])\)。

\begin{enumerate}
\def\labelenumi{\arabic{enumi})}
\setcounter{enumi}{6}
\tightlist
\item
  设 X 是一个可解圈拓扑空间，设 G 是一个在 X 上逆紧作用的离散群；记 \(m\colon G \times X \to X\) 为 G 的作用，\(\mathrm{pr}_2\colon G \times X \to X\) 为第二投影，\(f\colon X \to X / G\) 为典范满射。
\end{enumerate}

\begin{enumerate}
\def\labelenumi{\alph{enumi})}
\item
  证明群胚态射 \(\varpi(f)\) 是满射。
\item
  设 H 是一个群胚，设 \(\varphi\colon\varpi(X)\to H\) 是群胚态射，满足 \(\varphi\circ\varpi(m)=\varphi\circ\varpi(\mathrm{pr}_{2})\)。
\end{enumerate}

设 \(c_{1}\) 和 \(c_{2}\) 是 X 中的道路，使得道路 \(f \circ c_{1}\) 和 \(f \circ c_{2}\) 在 X/G 中严格同伦。证明 \(\varphi([c_{1}]) = \varphi([c_{2}])\)。

\begin{enumerate}
\def\labelenumi{\alph{enumi})}
\setcounter{enumi}{2}
\tightlist
\item
  证明存在唯一的群胚态射 \(\varphi'\)，从 \(\varpi(\mathrm{X/G})\) 到 H，使得 \(\varphi = \varphi' \circ \varpi(f)\)。由此得到 IV 第 403 页定理 3 的一个新证明。
\end{enumerate}

